\section{从数据矩阵到 PCA}
\label{sec:03-Linear-Algebra/07-SVD-and-Data-Applications/04}

手上一张表：\(N\) 行样本，\(d=50\) 列特征。你想画一张二维散点图给同事看，看看样本有没有分成几堆。二维意味着只能留两个坐标，于是很自然地想：挑两列最重要的，剩下四十八列扔掉。

问题是``重要''这件事跟你怎么记录数据有关。若点云主要沿一条斜线铺开，横轴和纵轴单独看都平平无奇，把两列各自的方差排个序也选不出这条斜线——它根本不在坐标轴上。删列这个动作被原始坐标系绑死了。

换个问法就松开了：能不能重新挑一组正交方向，让第一个方向尽量铺开点云，第二个方向在垂直于它的范围内继续尽量铺开，依此类推？这就是主成分分析（PCA）。

\begin{center}
\begin{tikzpicture}[x=1cm,y=1cm,font=\small,>=stealth]
  \draw[black!35,->] (-2.9,0) -- (2.9,0) node[right,black!50] {特征 1};
  \draw[black!35,->] (0,-1.9) -- (0,1.9) node[above,black!50] {特征 2};
  \begin{scope}[rotate=28]
    \draw[black!45,dashed] (0,0) ellipse [x radius=2.25,y radius=0.78];
    \fill (-2.00,0.15) circle (1.6pt);
    \fill (-1.62,-0.32) circle (1.6pt);
    \fill (-1.25,0.36) circle (1.6pt);
    \fill (-0.92,-0.14) circle (1.6pt);
    \fill (-0.55,0.46) circle (1.6pt);
    \fill (-0.20,-0.36) circle (1.6pt);
    \fill (0.12,0.20) circle (1.6pt);
    \fill (0.42,-0.47) circle (1.6pt);
    \fill (0.80,0.31) circle (1.6pt);
    \fill (1.12,-0.21) circle (1.6pt);
    \fill (1.48,0.41) circle (1.6pt);
    \fill (1.82,-0.30) circle (1.6pt);
    \fill (2.08,0.11) circle (1.6pt);
    \fill (-1.40,0.55) circle (1.6pt);
    \draw[->,very thick,blue!70!black] (0,0) -- (2.25,0);
    \node[blue!70!black] at (2.75,0) {\(v_1\)};
    \draw[->,very thick,blue!70!black] (0,0) -- (0,0.78);
    \node[blue!70!black] at (0,1.18) {\(v_2\)};
  \end{scope}
\end{tikzpicture}

\emph{灰色是记录数据时用的坐标轴，两条都不对齐点云的走向。蓝色是 PCA 选出的方向：\(v_1\) 沿点云铺得最开，\(v_2\) 垂直于它。虚线椭圆是等密度线的示意，它的两条半轴长度之比就是两个方向上标准差之比。}
\end{center}

\subsection{中心化不是预处理小节，是这套方法的前提}

在动手求方向之前必须先减去均值。令 \(\mu\) 为各列的均值向量，把每个样本换成 \(x_i-\mu\)。跳过这一步，接下来的所有公式仍然算得出数，只是算的东西变了：它们会找出数据相对于坐标原点的主要伸展方向，而不是相对于数据自身中心的变化方向。

差别在数据离原点较远时非常大。设想一批人的``身高（厘米）''与``体重（千克）''，点云中心在 \((170,65)\) 附近，各方向的展开幅度只有十几。不中心化时，第一个方向几乎就是从原点指向 \((170,65)\) 的那条射线——它描述的是``这批人平均有多大''，与个体之间的差异几乎无关。第二个方向才勉强开始描述你真正关心的变化。

\begin{center}
\begin{tikzpicture}[x=1cm,y=1cm,font=\footnotesize,>=stealth]
  \begin{scope}
    \draw[black!45,->] (-0.4,0) -- (3.9,0);
    \draw[black!45,->] (0,-0.4) -- (0,2.9);
    \node[below left,black!55] at (0,0) {\(O\)};
    \fill (1.371,1.037) circle (1.5pt);
    \fill (1.755,0.970) circle (1.5pt);
    \fill (1.797,1.398) circle (1.5pt);
    \fill (2.138,1.307) circle (1.5pt);
    \fill (2.220,1.666) circle (1.5pt);
    \fill (2.603,1.402) circle (1.5pt);
    \fill (2.680,1.769) circle (1.5pt);
    \fill (2.996,1.721) circle (1.5pt);
    \fill (2.060,1.642) circle (1.5pt);
    \fill (1.675,0.808) circle (1.5pt);
    \fill (2.483,1.910) circle (1.5pt);
    \fill (2.133,1.015) circle (1.5pt);
    \draw[->,very thick,red!70!black] (0,0) -- (2.75,1.75);
    \node[align=center,red!70!black] at (1.00,2.40) {未中心化时\\第一方向指向均值};
    \node at (1.75,-1.0) {原始坐标下的点云};
  \end{scope}
  \begin{scope}[shift={(7.4,1.4)}]
    \draw[black!45,->] (-1.6,0) -- (1.7,0);
    \draw[black!45,->] (0,-1.4) -- (0,1.4);
    \node[below right,black!55] at (0,0) {\(O\)};
    \fill (-0.829,-0.363) circle (1.5pt);
    \fill (-0.445,-0.430) circle (1.5pt);
    \fill (-0.403,-0.002) circle (1.5pt);
    \fill (-0.062,-0.093) circle (1.5pt);
    \fill (0.020,0.266) circle (1.5pt);
    \fill (0.403,0.002) circle (1.5pt);
    \fill (0.480,0.369) circle (1.5pt);
    \fill (0.796,0.321) circle (1.5pt);
    \fill (-0.140,0.242) circle (1.5pt);
    \fill (-0.525,-0.592) circle (1.5pt);
    \fill (0.283,0.510) circle (1.5pt);
    \fill (-0.067,-0.385) circle (1.5pt);
    \draw[->,very thick,blue!70!black] (0,0) -- (1.13,0.65);
    \draw[->,very thick,blue!70!black] (0,0) -- (-1.13,-0.65);
    \node[align=center,blue!70!black] at (-0.75,0.85) {中心化后\\第一方向沿离散度};
    \node at (0,-2.4) {减去均值后的同一批点};
  \end{scope}
\end{tikzpicture}

\emph{同一批点，同一套公式，差别只在有没有先减均值。左边找到的方向回答``这批数据整体在哪''，右边找到的方向回答``它们彼此差在哪''。忘记中心化是实践中最常见的一类错误，症状是第一主成分的解释方差比高得离谱（常常超过 \(99\%\)），而所有样本在它上面的得分几乎相同。}
\end{center}

下文的 \(X\in\R^{N\times d}\) 一律指已经中心化的数据矩阵。

\subsection{最大方差是一个特征值问题，而 SVD 直接把答案摆出来}

取一个单位方向 \(v\in\R^d\)。所有样本投到该方向上的坐标构成向量 \(Xv\)，平方和是 \(\norm{Xv}_2^2=v^{\trans}X^{\trans}Xv\)。第一主方向要解的就是

\[
\max_{\norm{v}=1}\;v^{\trans}X^{\trans}Xv .
\]

这是对称矩阵 \(X^{\trans}X\) 的 Rayleigh 商在单位球上的最大值问题。把 \(v\) 在 \(X^{\trans}X\) 的特征基下展开成 \(\sum_ic_iw_i\)，商就成了以 \(c_i^2\) 为权的特征值加权平均，最大值只能由最大特征值取到，取到它的正是对应的特征向量。样本协方差记为 \(S=X^{\trans}X/(N-1)\) 时，主方向就是 \(S\) 的特征向量，第 \(i\) 个方向上的样本方差是 \(\lambda_i(S)\)。

到这里还不需要 SVD，但对 \(X\) 直接做 SVD 可以省掉构造 \(X^{\trans}X\) 这一步——上一节刚说过这个动作会把条件数平方。写 \(X=U\Sigma V^{\trans}\)，则 \(X^{\trans}X=V\Sigma^2V^{\trans}\)，所以 \(V\) 的列就是主方向（也叫载荷方向），样本在这些方向上的新坐标是

\[
Z=XV=U\Sigma ,
\]

称为得分。第 \(i\) 个主成分的样本方差是 \(\sigma_i^2/(N-1)\)。载荷住在特征空间、每列长度为 \(d\)，得分按样本排列、每列长度为 \(N\)，两者维度都不同，混用会得到看起来能跑的错误代码。

\subsection{最大方差与最小重构误差是同一句话}

保留前 \(k\) 个方向时，\(Z_k=XV_k\)，重构 \(\widehat X=Z_kV_k^{\trans}=XV_kV_k^{\trans}\)，也就是把每个样本正交投影到前 \(k\) 维主子空间上。

为什么这个子空间是最好的？给定任何一个 \(k\) 维子空间，逐样本作正交投影已经是该子空间内的最优重构，这是投影定理。剩下的问题是在所有 \(k\) 维子空间中挑一个，而 \(\widehat X=XV_kV_k^{\trans}\) 的秩不超过 \(k\)，上一节的 Eckart--Young 定理正好回答：使 \(\norm{X-\widehat X}_F\) 最小的秩 \(k\) 矩阵就是截断 SVD，对应的子空间就是前 \(k\) 个右奇异向量张成的那个。``让投影方差最大''和``让重构误差最小''因此不是两个准则，而是同一个勾股关系的两端——总平方长度固定，投影分量变大等价于残差分量变小。

举个能一眼看穿的例子。三个已中心化的样本 \((2,1)\)、\((0,0)\)、\((-2,-1)\) 全部落在一条过原点的直线上，\(X^{\trans}X\) 的特征值是 \(10\) 和 \(0\)，第一主方向 \(v_1=(2,1)^{\trans}/\sqrt5\)，第二个方向上的得分全为零。一个主成分就能无误差重建整个数据集。给每个点加一点垂直于直线的抖动，第二个奇异值从零变成一个小正数，秩一重构就成了``把点投回最合适的那条直线''。

\subsection{解释方差比说的比你想的少}

常用的``解释方差比''是 \(\sigma_i^2/\sum_j\sigma_j^2\)。这个名字诱导过度解读。它只说欧氏平方变差中有多少落在该方向上，不说因果、不说业务重要性、也不说对下游预测有多大贡献——上一节关于``能量比例不是信息比例''的提醒在这里原样成立。一个方差很小的方向完全可能是分类边界所在。

特征的单位会直接改变结果。身高以米计、收入以元计时，收入那一列的数值尺度能压过其余所有特征，第一主成分基本就是收入。是否先标准化到单位方差没有机械答案：量纲不可比时标准化通常有帮助；绝对尺度本身有意义时（比如同一批传感器的多个通道），标准化反而会把真实的强弱差异抹平。预处理是建模假设的一部分，写报告时必须交代。

还有几条边界值得记住。PCA 找的是线性子空间，数据若躺在一个弯曲的流形上，投影会把结构折叠掉。主方向的正负号没有意义，相近或重复的特征值会让单个方向在重采样中乱转，但对应的子空间往往仍然稳定——这正是第一节关于奇异向量唯一性那条边界在数据里的样子。均值和平方误差对异常值都很敏感，一个离群点可以把第一主方向整个拽过去。样本量与维数相当时，样本协方差的谱本身就是失真的，\hyperref[sec:11-Mathematical-Statistics/09-Random-Matrix-Theory/03]{``谱偏差怎样影响 PCA 与随机投影''}给出了失真的具体形状。最后是流程上的：中心与主方向必须在训练集上确定，测试数据套用同一套变换，重新拟合一遍再比较等于把测试集的信息漏进了模型。\hyperref[sec:13-Machine-Learning/04-Dimensionality-Reduction/01]{``PCA 从最大方差到最小重构误差''}会把这些实践问题接着展开。

至此本章的四条线合上了：同一个 \(A=U\Sigma V^{\trans}\)，看作映射时给出主轴与四个基本子空间，看作方程时给出伪逆与条件数，看作近似时给出最优截断，看作数据时给出主成分。它们共用的那句话是，任何矩阵都可以读成``先转到一组正交坐标、沿各轴独立伸缩、再转到另一组正交坐标''。

奇异向量是从矩阵里算出来的，这既是 SVD 的力量，也意味着每换一个矩阵就得重算一次。下一章问相反的问题：有没有一组事先固定的正交基，能让一整类矩阵同时变成对角？把标量域扩到复数之后，\hyperref[ch:024]{``复数、谱与 Fourier 结构''}给出的答案是 Fourier 基，而``事先固定''这一点正是快速算法得以存在的原因。
